\section{Metodología}
\label{sec:metodologia}

El presente proyecto se desarrollará bajo un enfoque metodológico que se fundamenta en el Manual de Presentación de Proyectos de la Facultad de Ingeniería en Ciencias de la Computación y Telecomunicaciones \textcite{calderon2013manual}, combinado con Scrum como marco de gestión ágil y la metodología UWE (UML-based Web Engineering) para el modelado y desarrollo del sistema.

\subsection{Scrum para la gestión del proyecto}

Se empleará Scrum para organizar y controlar el avance del proyecto mediante iteraciones denominadas sprints. El equipo de trabajo está conformado por tres desarrolladores, que asumen los roles de Scrum Master y equipo de desarrollo, mientras que el rol de Product Owner es omitido y asumido por el Scrum Master guiado por los casos de uso planteados y los requerimientos espeficados. Se llevarán a cabo las ceremonias habituales: Sprint Planning al inicio de cada iteración, Daily Standup para la sincronización diaria, Sprint Review para la presentación de los incrementos y Sprint Retrospective para la mejora continua. El backlog del producto se construirá a partir de los casos de uso definidos y se priorizará en función de las necesidades de la facultad.

\subsection{UWE para el modelado y desarrollo}

Para el análisis, diseño e implementación del sistema se utilizará la metodología UWE (UML-based Web Engineering) propuesta por \textcite{koch2001uwe}, la cual extiende UML con estereotipos específicos para aplicaciones web. UWE define cuatro modelos principales:

\begin{enumerate}
    \item \textbf{Modelo de contenido:} define la estructura de la información del sistema mediante diagramas de clases UML, identificando entidades como usuario, manuscrito, evaluación, pago, entre otras.
    \item \textbf{Modelo de navegación:} especifica las rutas de navegación y los accesos a la información según el perfil del usuario (vicedecano, director, tutor, tribunal, estudiante).
    \item \textbf{Modelo de presentación:} describe la interfaz de usuario y la disposición de los elementos visuales.
    \item \textbf{Modelo de procesos:} representa los flujos de trabajo académicos (registro, asignación, revisión, aprobación, publicación).
\end{enumerate}

\subsection{Correspondencia con los objetivos específicos}

La metodología propuesta permite alcanzar cada objetivo específico de la siguiente manera:

\begin{enumerate}
    \item El diagnóstico de los procesos actuales se realizará mediante entrevistas con autoridades, docentes y estudiantes de la FICH, así como la revisión de la documentación institucional.
    \item Los requisitos funcionales y técnicos se determinarán a partir del análisis de los casos de uso aprobados y se documentarán mediante diagramas UML.
    \item El diseño de los flujos de trabajo se modelará con UWE, representando los procesos de registro, asignación, evaluación y publicación.
    \item El desarrollo de los módulos se ejecutará en sprints Scrum, implementando incrementalmente cada caso de uso.
    \item La evaluación del impacto se realizará mediante indicadores cuantitativos (tiempos de gestión, número de proyectos finalizados) y cualitativos (encuestas de satisfacción a usuarios).
\end{enumerate}

\subsection{Recursos y herramientas}

Se emplearán herramientas de código abierto: el lenguaje PHP con el framework Laravel version 12 para el backend usando inertia con ReactJS para el frontend, una base de datos relacional PostgreSQL, y Git para el control de versiones. La gestión del proyecto se apoyará en tableros Kanban digitales para el seguimiento del backlog usando la herramienta Taiga.

\endinput
